\chapter{动量与冲量}

牛顿第二定律告诉我们：力决定加速度；而微积分进一步告诉我们，真正适合在瞬时变化、短时相互作用和多体系统中使用的量，是\emph{动量}而不是速度本身。动量把质量与速度统一起来，使“作用一段时间的力”可以自然地用积分表达；它又在系统观点下导出守恒定律，从而成为碰撞、爆炸、反冲、火箭推进等问题的核心语言。

本章将用微积分把下列链条连贯起来：
\[
\vb F=\dv{\vb p}{t}
\quad\Longrightarrow\quad
\vb I=\int_{t_1}^{t_2}\vb F\,\dd t=\Delta \vb p
\quad\Longrightarrow\quad
\text{封闭系统总动量守恒}.
\]
这条链条既能解释短暂碰撞，也能处理变质量系统中的持续喷气过程。与前一章相比，本章的重点从“受力后如何加速”转向“相互作用前后整体如何变化”。

\section{动量的定义}

\subsection{质量与速度的乘积}

\begin{definition}
质点的\emph{动量}定义为
\[
\vb p=m\vb v,
\]
其中 $m$ 是质点质量，$\vb v$ 是瞬时速度。动量是一个与速度方向相同的矢量，其单位是 $\mathrm{kg\cdot m/s}$。
\end{definition}

如果质量恒定，则动量与速度成正比；质量越大、速度越快，物体就越“难以被改变运动状态”。高中阶段常把这种“运动惯性”说得较为直观，而在微积分语言中，它最自然的刻画正是 $\vb p=m\vb v$。

\begin{remark}
动量是矢量，不仅有大小，还有方向。两个物体即使速率相同，只要方向不同，动量也不同。特别地，在一维问题中常把向右规定为正方向，于是动量可以用带正负号的标量表示；但在二维、三维问题中必须保留矢量形式。
\end{remark}

把动量写成分量形式，若
\[
\vb v=v_x\vb e_x+v_y\vb e_y+v_z\vb e_z,
\]
则
\[
\vb p=mv_x\vb e_x+mv_y\vb e_y+mv_z\vb e_z.
\]
因此动量在各坐标方向上的分量分别为 $p_x=mv_x$、$p_y=mv_y$、$p_z=mv_z$。这意味着在外力某一方向分量为零时，该方向动量可能守恒。

\begin{example}
质量为 $0.20\,\mathrm{kg}$ 的小球，以速率 $5\,\mathrm{m/s}$ 沿东偏北 $37^\circ$ 方向运动。求其动量矢量。

解：取东向为 $x$ 轴、北向为 $y$ 轴，则
\[
\vb v=5\cos 37^\circ\,\vb e_x+5\sin 37^\circ\,\vb e_y.
\]
故
\[
\vb p=0.20\vb v=(1.0\cos 37^\circ)\vb e_x+(1.0\sin 37^\circ)\vb e_y.
\]
近似取 $\cos37^\circ\approx0.8$，$\sin37^\circ\approx0.6$，得
\[
\vb p\approx0.80\,\vb e_x+0.60\,\vb e_y\quad (\mathrm{kg\cdot m/s}).
\]
其大小为 $1.0\,\mathrm{kg\cdot m/s}$，方向与速度相同。
\end{example}

\subsection{为什么动量比速度更基本}

牛顿第二定律在最一般的形式中并不是 $\vb F=m\vb a$，而是
\[
\vb F=\dv{\vb p}{t}.
\]
当质量恒定时，
\[
\dv{\vb p}{t}=\dv{(m\vb{v})}{t}=m\dv{\vb v}{t}=m\vb a,
\]
这才退化成熟悉的 $\vb F=m\vb a$。因此，动量是牛顿动力学中更基本的状态量；尤其在火箭这类变质量问题里，必须回到 $\vb F=\dv{\vb p}{t}$ 才能正确处理。

\begin{remark}
从微积分角度看，速度是位置对时间的一阶导数，而动量是“质量乘以速度”得到的物理量。它一方面直接与力的时间积分相联系，另一方面在系统内部相互作用中能发生漂亮的抵消，因此在多体问题里特别重要。
\end{remark}

\subsection{质点系的总动量}

\begin{definition}
由若干个质点组成的系统，其\emph{总动量}定义为各质点动量的矢量和：
\[
\vb P=\sum_{i=1}^N \vb p_i=\sum_{i=1}^N m_i\vb v_i.
\]
若系统连续分布，也可写成积分形式
\[
\vb P=\int \vb v\,\dd m.
\]
\end{definition}

总动量强调“整体视角”。在两体、三体乃至连续介质中，单个质点的速度可能都在变化，但只要系统满足适当条件，总动量仍可能保持不变。

\section{冲量与动量定理}

\subsection{从微分到积分}

\begin{law}[动量定理]
若质点在时间区间 $[t_1,t_2]$ 内所受合外力为 $\vb F(t)$，则
\[
\vb F=\dv{\vb p}{t}
\]
在该区间上积分得到
\[
\int_{t_1}^{t_2}\vb F(t)\,\dd t=\vb p(t_2)-\vb p(t_1)=\Delta\vb p.
\]
力对时间的积分称为\emph{冲量}，记作
\[
\vb I\coloneqq \int_{t_1}^{t_2}\vb F(t)\,\dd t.
\]
因此
\[
\vb I=\Delta\vb p.
\]
\end{law}

这个定律说明：力不只是“有没有”，还要看“作用了多久”。一个很大的力如果只作用极短时间，其效果未必比一个较小但持续更久的力更显著。冲量正是把“力的大小”和“持续时间”统一起来的量。

\begin{remark}
若力恒定，则 $\vb I=\vb F\Delta t$；若力随时间变化，就必须真正进行积分。图像上，冲量等于 $F$--$t$ 图像与时间轴围成的有向面积；方向由矢量方向或一维符号决定。
\end{remark}

\subsection{平均力与缓冲过程}

在工程和生活中，人们常关心“同样的动量变化，怎样减小峰值作用力”。若某过程的总冲量固定为 $\vb I$，持续时间为 $\Delta t$，则可定义平均力
\[
\avg{\vb F}=\frac{1}{\Delta t}\int_{t_1}^{t_2}\vb F\,\dd t=\frac{\vb I}{\Delta t}.
\]
这说明：在动量改变量 $\Delta\vb p$ 已确定时，延长受力时间可以减小平均力。这就是安全气囊、缓冲垫、沙坑、海绵包装能够“减震”的动力学本质。

严格地说，平均力并不等于最大力。真实碰撞中，$F(t)$ 往往先迅速增大后减小；但只要把受力时间拉长，在同样冲量下，曲线下的面积不变而峰值通常会下降。微积分把这种经验判断转化成了清楚的面积比较问题。

\begin{figure}[htbp]
\centering
\begin{tikzpicture}
\begin{axis}[
  width=0.75\textwidth, height=0.5\textwidth,
  xlabel={时间 $t$}, ylabel={力 $F$},
  axis lines=middle,
  grid=major, grid style={gray!30},
  thick,
  xmin=0, xmax=5.2, ymin=0, ymax=3.5,
  xtick={0.8,4.2},
  xticklabels={$t_1$,$t_2$},
]
\addplot[name path=force, blue, domain=0.8:4.2, samples=120] {3*exp(-1.2*(x-2.5)^2)};
\addplot[name path=axisline, draw=none, domain=0.8:4.2] {0};
\addplot[blue!20] fill between[of=force and axisline];
\end{axis}
\end{tikzpicture}
\caption{变力冲量等于 $F$--$t$ 图像下的阴影面积：峰值很高但作用时间很短}
\end{figure}

\subsection{变力冲量的积分计算}

当作用力明显随时间变化时，最可靠的方法就是积分。例如，若一维情况下
\[
F(t)=kt\qquad (0\le t\le T),
\]
则冲量为
\[
I=\int_0^T kt\,\dd t=\frac12 kT^2.
\]
如果力先正后负，就要按定积分处理正负面积，而不是简单把绝对值相加。

\begin{example}
某物体在 $0\le t\le \tau$ 内受一维变力
\[
F(t)=F_0\sin\frac{\pi t}{\tau},
\]
初动量为 $p_0$。求末动量。

解：由动量定理，
\[
I=\int_0^{\tau}F_0\sin\frac{\pi t}{\tau}\,\dd t
=\frac{F_0\tau}{\pi}\int_0^{\pi}\sin u\,\dd u
=\frac{2F_0\tau}{\pi}.
\]
故
\[
p(\tau)=p_0+\frac{2F_0\tau}{\pi}.
\]
若物体质量为常量 $m$，则其速度变化量为
\[
\Delta v=\frac{1}{m}\cdot \frac{2F_0\tau}{\pi}.
\]
这个结果提醒我们：碰撞或推力脉冲的效果，往往取决于力函数积分后的总面积，而不只取决于峰值 $F_0$。
\end{example}

\subsection{动量定理的分量形式}

在二维或三维问题中，动量定理可写成
\[
\int_{t_1}^{t_2}\vb F\,\dd t=\Delta \vb p,
\]
也可按坐标方向分别写为
\[
\int_{t_1}^{t_2}F_x\,\dd t=\Delta p_x,
\qquad
\int_{t_1}^{t_2}F_y\,\dd t=\Delta p_y.
\]
这在斜碰、斜抛爆炸、带约束碰撞等题目里很有用：某一方向如果外冲量为零，则该方向动量单独守恒。

\begin{remark}
短时碰撞中，重力等通常持续存在，但其冲量 $mg\Delta t$ 往往远小于碰撞力产生的冲量，可在近似中忽略。这就是“碰撞瞬间系统动量近似守恒”的微积分依据：不是外力为零，而是\emph{外冲量可忽略}。
\end{remark}

\section{动量守恒定律}

\subsection{由牛顿第三定律导出}

考虑两个质点 $1,2$ 组成的系统。设它们之间的相互作用力分别为 $\vb F_{12}$ 和 $\vb F_{21}$，外力分别为 $\vb F_1^{\text{ext}}$、$\vb F_2^{\text{ext}}$。由动量定理的微分形式，
\[
\dv{\vb p_1}{t}=\vb F_1^{\text{ext}}+\vb F_{21},
\qquad
\dv{\vb p_2}{t}=\vb F_2^{\text{ext}}+\vb F_{12}.
\]
两式相加得
\[
\dv{(\vb{p}_1+\vb{p}_2)}{t}=\vb F_1^{\text{ext}}+\vb F_2^{\text{ext}}+(\vb F_{12}+\vb F_{21}).
\]
若满足牛顿第三定律，
\[
\vb F_{12}=-\vb F_{21},
\]
则内力项相互抵消，于是
\[
\dv{\vb P}{t}=\vb F_{\text{ext}},
\]
其中 $\vb P=\vb p_1+\vb p_2$ 是系统总动量，$\vb F_{\text{ext}}$ 是系统所受外力的矢量和。

对 $N$ 个质点的系统，这个结论同样成立：成对内力在总和中抵消，故总动量的变化只由外力决定。

\begin{law}[动量守恒定律]
若一个系统所受合外力为零，或者在所研究时间内外冲量可以忽略，则系统总动量保持不变，即
\[
\dv{\vb P}{t}=0
\quad\Longrightarrow\quad
\vb P=\text{常量}.
\]
等价地，
\[
\vb P_{\text{初}}=\vb P_{\text{末}}.
\]
\end{law}

对 $N$ 个质点组成的系统，把第 $i$ 个质点所受外力记为 $\vb F_i^{\text{ext}}$，所受其他质点内力之和记为 $\sum_{j\ne i}\vb F_{ij}$，则
\[
\dv{\vb p_i}{t}=\vb F_i^{\text{ext}}+\sum_{j\ne i}\vb F_{ij}.
\]
对全部 $i$ 求和，得到
\[
\dv{}{t}\!\left(\sum_{i=1}^N\vb{p}_i\right)=\sum_{i=1}^N\vb F_i^{\text{ext}}+\sum_{i=1}^N\sum_{j\ne i}\vb F_{ij}.
\]
如果内力满足成对反向，则双重求和中的各项两两抵消，只剩外力总和。这一推导说明，总动量守恒并不是“巧合公式”，而是由微分方程对整个系统求和后出现的结构性结果。

\subsection{应用条件}

动量守恒并不是“任何时候都能直接套公式”。要使用它，必须先检查系统边界与外力条件：
\begin{enumerate}
  \item 是否选择了合适的研究系统；
  \item 系统所受合外力是否为零；
  \item 若外力不为零，研究过程是否足够短，以至于外冲量可忽略；
  \item 若只在某一方向外力为零，是否只能在该方向使用动量守恒。
\end{enumerate}

例如，人在冰面上推箱子，若把“人+箱子”作为系统，水平外力近似为零，则总动量守恒；若只把箱子作为系统，则人的推力属于外力，箱子动量不守恒。

\begin{example}
两名滑冰者甲、乙静止于光滑冰面上，质量分别为 $60\,\mathrm{kg}$ 和 $40\,\mathrm{kg}$。甲向东推乙后，乙向东速度为 $3\,\mathrm{m/s}$。求甲的速度。

解：取“甲+乙”为系统。水平方向外力近似为零，初总动量为零，故末总动量仍为零：
\[
60v_\text{甲}+40\times 3=0.
\]
解得
\[
v_\text{甲}=-2\,\mathrm{m/s}.
\]
负号表示甲向西运动。这个结果说明：动量守恒体现的是整体平衡，而不是各自速度相等。
\end{example}

\begin{remark}
动量守恒是矢量守恒。二维问题中常拆成两个独立方程：
\[
P_x=\text{常量},\qquad P_y=\text{常量}.
\]
这往往比只写一个矢量式更便于计算。
\end{remark}

\section{碰撞}

碰撞是动量方法最典型的应用。所谓碰撞，并不局限于硬球相撞，也包括子弹射入木块、两车挂接、球拍击球、爆炸分离等短时强相互作用过程。其共同特征是：接触时间很短、相互作用力很大，因而最适合使用冲量与动量。

\subsection{碰撞过程的微积分描述}

设两物体在 $t_1$ 到 $t_2$ 的短时间内发生相互作用。对其中一个物体有
\[
\Delta \vb p_1=\int_{t_1}^{t_2}(\vb F_{21}+\vb F_1^{\text{ext}})\,\dd t,
\]
对另一个物体有
\[
\Delta \vb p_2=\int_{t_1}^{t_2}(\vb F_{12}+\vb F_2^{\text{ext}})\,\dd t.
\]
若外冲量可以忽略，则
\[
\Delta \vb p_1+\Delta \vb p_2
=\int_{t_1}^{t_2}(\vb F_{12}+\vb F_{21})\,\dd t=0.
\]
因此总动量守恒。

这段推导揭示了碰撞分析的本质：不是把碰撞看成“速度瞬间神秘改变”，而是把它看成一个极短时间内的强变力积分过程。速度的突变，是大力在短时间内积分后的结果。

\begin{figure}[H]
\centering
\begin{tikzpicture}[scale=1.0,>=Stealth]
  \draw[thick] (0,0) -- (7.6,0);
  \filldraw[fill=blue!20] (1.0,0) rectangle (2.0,0.8);
  \filldraw[fill=red!20] (5.2,0) rectangle (6.4,1.0);
  \draw[->,thick,blue!70!black] (0.2,0.4) -- (0.95,0.4) node[midway,above] {$u_1$};
  \draw[->,thick,red!70!black] (4.2,0.5) -- (5.15,0.5) node[midway,above] {$u_2$};
  \node at (1.5,-0.35) {$m_1$};
  \node at (5.8,-0.35) {$m_2$};
  \draw[->,thick,black] (2.4,0.4) -- (4.8,0.4) node[midway,above] {碰撞区间};
\end{tikzpicture}
\caption{一维碰撞示意图：碰撞前后可用动量定理连接}
\end{figure}

\subsection{弹性碰撞与非弹性碰撞}

\begin{definition}
碰撞前后若系统总动能保持不变，称为\emph{弹性碰撞}；若总动能减少，称为\emph{非弹性碰撞}。其中碰后两物体粘在一起共同运动的情形，称为\emph{完全非弹性碰撞}。
\end{definition}

无论弹性还是非弹性，只要外冲量可忽略，总动量都守恒。差别在于动能是否守恒。换言之，\emph{动量守恒比机械能守恒更普遍}。

\begin{theorem}[一维弹性碰撞速度公式]
设两物体质量分别为 $m_1,m_2$，碰前速度为 $u_1,u_2$，碰后速度为 $v_1,v_2$。若碰撞为一维弹性碰撞，则
\[
 v_1=\frac{m_1-m_2}{m_1+m_2}u_1+\frac{2m_2}{m_1+m_2}u_2,
\]
\[
 v_2=\frac{2m_1}{m_1+m_2}u_1+\frac{m_2-m_1}{m_1+m_2}u_2.
\]
\end{theorem}

\begin{proof}
由动量守恒得
\[
 m_1u_1+m_2u_2=m_1v_1+m_2v_2.
\]
由动能守恒并整理，可得到相对速度关系
\[
 u_1-u_2=-(v_1-v_2).
\]
于是
\[
 v_1-v_2=-(u_1-u_2).
\]
将它与动量方程联立求解即可得到上式。可见弹性碰撞的求解，本质上是一个线性方程组问题，而“相对速度反向”则是机械能守恒压缩后的简洁表达。
\end{proof}

对于一维弹性碰撞，设碰前速度为 $u_1,u_2$，碰后速度为 $v_1,v_2$，则有
\[
m_1u_1+m_2u_2=m_1v_1+m_2v_2,
\]
以及
\[
\frac12 m_1u_1^2+\frac12 m_2u_2^2
=\frac12 m_1v_1^2+\frac12 m_2v_2^2.
\]
联立可推出相对速度关系
\[
u_1-u_2=-(v_1-v_2).
\]
这表示：弹性碰撞中，分离速度等于接近速度，方向相反。

\begin{figure}[htbp]
\centering
\begin{minipage}{0.48\textwidth}
\centering
\begin{tikzpicture}
\begin{axis}[
  width=\linewidth, height=0.7\linewidth,
  xlabel={时间 $t$}, ylabel={速度 $v$},
  title={弹性碰撞},
  axis lines=middle,
  grid=major, grid style={gray!30},
  thick,
  xmin=0, xmax=4.2, ymin=-0.2, ymax=1.8,
  legend style={at={(0.98,0.98)}, anchor=north east, draw=none, fill=white},
]
\addplot[blue] coordinates {(0,1.5) (2,1.5)};
\addplot[blue] coordinates {(2,0.2) (4,0.2)};
\addplot[red] coordinates {(0,0) (2,0)};
\addplot[red] coordinates {(2,1.3) (4,1.3)};
\legend{$v_1(t)$,$v_2(t)$}
\end{axis}
\end{tikzpicture}
\end{minipage}
\hfill
\begin{minipage}{0.48\textwidth}
\centering
\begin{tikzpicture}
\begin{axis}[
  width=\linewidth, height=0.7\linewidth,
  xlabel={时间 $t$}, ylabel={速度 $v$},
  title={完全非弹性碰撞},
  axis lines=middle,
  grid=major, grid style={gray!30},
  thick,
  xmin=0, xmax=4.2, ymin=-0.2, ymax=1.8,
  legend style={at={(0.98,0.98)}, anchor=north east, draw=none, fill=white},
]
\addplot[blue] coordinates {(0,1.5) (2,1.5)};
\addplot[blue] coordinates {(2,0.5) (4,0.5)};
\addplot[red] coordinates {(0,0) (2,0)};
\addplot[red] coordinates {(2,0.5) (4,0.5)};
\legend{$v_1(t)$,$v_2(t)$}
\end{axis}
\end{tikzpicture}
\end{minipage}
\caption{碰撞前后速度的阶跃变化：弹性碰撞交换更多动能，而完全非弹性碰撞在碰后共速前进}
\end{figure}

\begin{theorem}[一维完全非弹性碰撞]
若质量为 $m_1,m_2$ 的两物体碰后粘连，以共同速度 $v$ 运动，则
\[
v=\frac{m_1u_1+m_2u_2}{m_1+m_2}.
\]
碰撞损失的机械能为
\[
\Delta E=\frac12 m_1u_1^2+\frac12 m_2u_2^2-\frac12 (m_1+m_2)v^2\ge 0.
\]
\end{theorem}

\begin{example}
质量分别为 $m$ 和 $2m$ 的两小球在光滑直线上碰撞。碰前前者速度为 $u$，后者静止。若碰撞为弹性碰撞，求碰后速度。

解：设碰后速度分别为 $v_1,v_2$。由动量守恒与动能守恒，
\[
mu=mv_1+2mv_2,
\qquad
\frac12 mu^2=\frac12 mv_1^2+m v_2^2.
\]
约去公共因子后得
\[
u=v_1+2v_2,
\qquad
u^2=v_1^2+2v_2^2.
\]
也可用相对速度关系
\[
u-0=-(v_1-v_2),
\]
即
\[
v_2-v_1=u.
\]
与 $u=v_1+2v_2$ 联立，解得
\[
v_1=-\frac{u}{3},\qquad v_2=\frac{2u}{3}.
\]
因此轻球反弹，重球向前运动。这与经验一致：轻球撞上较重静止物体后常会弹回。
\end{example}

\begin{example}
质量为 $m$ 的子弹以速度 $v_0$ 水平射入静止木块，木块质量为 $M$，子弹未穿出并留在木块中。求碰后的共同速度，并说明为何此过程不守恒机械能。

解：碰撞极短，可忽略外冲量，故动量守恒：
\[
mv_0=(M+m)v.
\]
所以
\[
v=\frac{m}{M+m}v_0.
\]
碰前后机械能分别为
\[
E_i=\frac12 mv_0^2,
\qquad
E_f=\frac12 (M+m)v^2=\frac12\frac{m^2}{M+m}v_0^2.
\]
显然 $E_f<E_i$。减少的机械能转化为内能、形变能、声能等。这正是完全非弹性碰撞的典型特征：动量守恒，机械能不守恒。
\end{example}

\begin{remark}
碰撞题常有两个层次：
\begin{enumerate}
  \item 先用短时过程中的动量关系求碰后瞬时速度；
  \item 再对碰后较长过程使用能量、运动学或牛顿定律分析。
\end{enumerate}
例如“子弹木块”问题中，碰撞后的整体还可能继续压缩弹簧或在粗糙面上滑行，此时动量守恒通常不再适用，但机械能或动能定理可能适用。
\end{remark}

\section{质心与质心运动}

\subsection{质心的定义}

\begin{definition}
设系统由若干质点组成，总质量为
\[
M=\sum_{i=1}^N m_i.
\]
其\emph{质心位置矢量}定义为
\[
\vb R=\frac{1}{M}\sum_{i=1}^N m_i\vb r_i.
\]
若质量连续分布，则
\[
\vb R=\frac{1}{M}\int \vb r\,\dd m.
\]
\end{definition}

若只研究一条直线上的分布，质心坐标就是
\[
X=\frac{1}{M}\int x\,\dd m.
\]
这与统计中的加权平均完全类似：质量越大的部分，对质心位置的“拉拽”越强。对于均匀细杆、薄板、圆环等连续体，质心积分常可转化为长度积分、面积积分或体积分。

\begin{figure}[H]
\centering
\begin{tikzpicture}[scale=1.0,>=Stealth]
  \draw[thick] (0,0) -- (7,0);
  \filldraw[fill=blue!20] (1.2,0) circle (0.12);
  \filldraw[fill=red!20] (5.5,0) circle (0.16);
  \filldraw[fill=black] (4.07,0) circle (0.08);
  \node[above] at (1.2,0.12) {$m_1$};
  \node[above] at (5.5,0.16) {$m_2$};
  \node[below] at (4.07,-0.02) {$R$};
  \draw[decorate,decoration={brace,amplitude=5pt}] (1.2,-0.35) -- (4.07,-0.35) node[midway,below=6pt] {$R-r_1$};
  \draw[decorate,decoration={brace,amplitude=5pt}] (4.07,-0.35) -- (5.5,-0.35) node[midway,below=6pt] {$r_2-R$};
\end{tikzpicture}
\caption{两质点系统的质心更靠近质量较大的一侧}
\end{figure}

质心可以理解为“质量分布的平均位置”。若系统总质量固定，则对时间求导得到
\[
\vb V_{\text{cm}}=\dv{\vb R}{t}=\frac{1}{M}\sum_{i=1}^N m_i\vb v_i,
\]
从而
\[
M\vb V_{\text{cm}}=\sum_{i=1}^N m_i\vb v_i=\vb P.
\]
这表明：\emph{系统总动量等于总质量乘以质心速度}。

\begin{theorem}[质心运动定理]
对于总质量恒定的质点系，质心满足
\[
M\dv[2]{\vb R}{t}=\vb F_{\text{ext}}.
\]
等价地，
\[
\dv{(M\vb{V}_{\text{cm}})}{t}=\vb F_{\text{ext}}.
\]
也就是说，系统质心的运动只由外力决定，内部相互作用不会改变质心运动规律。
\end{theorem}

\begin{proof}
由 $\vb P=M\vb V_{\text{cm}}$，再结合
\[
\dv{\vb P}{t}=\vb F_{\text{ext}},
\]
即可得到
\[
M\dv{\vb V_{\text{cm}}}{t}=M\dv[2]{\vb R}{t}=\vb F_{\text{ext}}.
\]
证明完毕。
\end{proof}

\begin{remark}
这个定理极其重要。烟花在空中爆炸时，碎片四散飞出，看起来非常复杂；但若忽略空气阻力，整个系统质心仍按原抛体轨道运动。复杂多体运动的“主线”，就是质心运动。
\end{remark}

\begin{example}
质量分别为 $m$ 和 $2m$ 的两个质点位于 $x$ 轴上，位置分别为 $x_1=0$、$x_2=3L$。求系统质心位置。

解：在一维中，质心坐标为
\[
X=\frac{mx_1+2mx_2}{3m}=\frac{0+2m\cdot 3L}{3m}=2L.
\]
可见质心更靠近质量较大的质点。这说明质心不是几何中点，而是按质量加权的平均位置。
\end{example}

\subsection{质心与守恒定律}

当系统所受合外力为零时，
\[
M\dv[2]{\vb R}{t}=0,
\]
故质心做匀速直线运动；若系统初始静止，则质心保持静止。这一结论常用于爆炸与反冲问题。

\begin{example}
一枚礼花弹在最高点爆炸成质量相等的两块碎片。若忽略空气阻力，则爆炸后两块碎片的质心如何运动？

解：把两碎片视为系统。爆炸属于内力作用，系统所受外力只有重力，总外力与爆炸前相同。因此质心继续做抛体运动。又因为爆炸发生在最高点，质心在该瞬间的竖直速度为零，之后按自由落体方式下降，同时保持原有的水平速度。无论碎片如何分离，质心轨迹都不受内部爆炸细节影响。
\end{example}

\section{火箭方程}

火箭推进的关键在于：火箭把部分燃料以高速向后喷出，从而使自身获得向前的动量。这个过程不是恒质量问题，因此不能直接使用 $\vb F=m\vb a$，而必须回到动量的基本形式。

\subsection{变质量系统的微元分析}

设某时刻火箭质量为 $m$，速度为 $v$。经过极短时间 $\dd t$ 后，火箭喷出一小部分燃气，质量减小为 $m+\dd m$，其中 $\dd m<0$。设喷气相对火箭的速度大小为 $u$，方向向后；则在地面参考系中，喷出气体速度约为 $v-u$。

若忽略外力，在时间 $\dd t$ 前后的总动量守恒：
\[
mv=(m+\dd m)(v+\dd v)+(-\dd m)(v-u).
\]
忽略二阶小量 $\dd m\,\dd v$，展开得
\[
mv=mv+m\dd v+v\dd m- v\dd m+u\dd m,
\]
于是
\[
m\dd v+u\dd m=0.
\]
由于 $\dd m<0$，可写成
\[
m\dd v=-u\dd m.
\]
两边积分得
\[
\int_{v_0}^{v}\dd v=-u\int_{m_0}^{m}\frac{1}{m}\,\dd m,
\]
即
\[
v-v_0=u\ln\frac{m_0}{m}.
\]

\begin{law}[齐奥尔科夫斯基火箭方程]
在忽略外力且喷气相对速度 $u$ 为常量时，火箭速度满足
\[
v=v_0+u\ln\frac{m_0}{m}.
\]
其中 $m_0$ 为初始质量，$m$ 为某时刻剩余质量。
\end{law}

\subsection{含外力的微分形式}

若还考虑外力 $\vb F_{\text{ext}}$，则矢量形式可写为
\[
\vb F_{\text{ext}}=\dv{(m\vb{v})}{t}-\vb u\,\dv{m}{t},
\]
这里 $\vb u$ 表示喷气相对火箭的速度矢量。对于竖直发射、忽略空气阻力的情形，若取向上为正，则有
\[
m\dv{v}{t}=-u\dv{m}{t}-mg.
\]
可见火箭加速并不只由“推力大不大”决定，还取决于当时剩余质量是否已经足够小。

\begin{remark}
火箭方程中的对数来自积分
\[
\int \frac{1}{m}\,\dd m=\ln m.
\]
这正是微积分在物理中的典型作用：把连续喷气分解成无数个“喷出一个极小质量元”的过程，再把这些微小速度增量累加起来。
\end{remark}

\begin{example}
一枚火箭在深空中飞行，可忽略外力。初始质量为 $5m$，燃料燃尽后剩余质量为 $m$，喷气相对速度恒为 $u$，初速度为零。求末速度。

解：由齐奥尔科夫斯基方程，
\[
v=u\ln\frac{5m}{m}=u\ln 5.
\]
因此末速度不与喷出燃料的总质量简单成正比，而与质量比的对数有关。若想进一步提高末速度，仅增加少量燃料并不总是高效，这也是多级火箭设计的数学背景之一。
\end{example}

\section*{本章小结}

本章围绕“力对时间的累积效应”和“系统内部作用的相互抵消”建立了动力学的第二条主线：
\begin{enumerate}
  \item 动量定义为 $\vb p=m\vb v$，是比速度更适合处理瞬时相互作用的量；
  \item 冲量是力对时间的积分，满足 $\vb I=\Delta\vb p$；
  \item 系统总动量变化只由外力决定，因此在封闭系统中总动量守恒；
  \item 碰撞问题的本质是短时强变力积分；
  \item 质心满足“总质量乘以质心加速度等于合外力”；
  \item 变质量系统必须使用动量基本形式，进而导出火箭方程。
\end{enumerate}
微积分并没有替代高中物理中的结论，而是揭示了这些结论为何成立、在什么条件下成立，以及在更复杂情形下该如何推广。

